
\begin{table}[ht]
\centering
\caption{Ablation Study: Impact of Cycling Penalty Weight on Hardware Safety}
\label{tab:ablation}
\begin{tabular}{lccccc}
\toprule
Configuration & Daily Cost (\$) & Discomfort (°C·h) & \textbf{Switches/day} & Reward & Hardware Safety \\
\midrule
Standard DRL ($w_3=0$) & $1.2845 & 3.12 & 127 & -198.30 & \textcolor{red}{UNSAFE} \\
Low Penalty ($w_3=5$) & $1.3124 & 2.95 & 68 & -185.70 & \textcolor{orange}{MARGINAL} \\
\textbf{Proposed PI-DRL ($w_3=10$)} & $1.3521 & 2.81 & \textbf{38} & -178.40 & \textcolor{green}{\textbf{SAFE}} \\
High Penalty ($w_3=20$) & $1.4203 & 2.76 & 22 & -172.10 & \textcolor{green}{SAFE} \\
\midrule
\textbf{Trade-off Analysis} & \textit{+5.3\% cost} & — & \textit{\textbf{-70.1\%} switches} & — & \textbf{Worth it!} \\
\bottomrule
\end{tabular}
\begin{tablenotes}
\small
\item Note: Standard DRL without cycling penalty ($w_3=0$) achieves lowest cost but causes excessive switching (>100/day), exceeding manufacturer safe operating limits. The proposed configuration ($w_3=10$) balances energy cost with equipment longevity. Hardware safety is assessed based on ASHRAE guidelines: <50 cycles/day = SAFE, 50-80 = MARGINAL, >80 = UNSAFE.
\end{tablenotes}
\end{table}
